\section{Introduction}%1 page
\label{sec:introduction}

% 注意多加引用

% 第一段：深度学习应用的强大能力使其活跃于各个领域...从CV到NLP，再到大模型、多模态...范式在逐渐演化

% 第二段：深度学习的计算量很大，使其适合在GPU或领域特定加速器上部署...为了提高性能，厂商提出了cuDNN、xx等计算库..然而，开发高性能库非常耗时，且可能跟不上快速迭代的网络形态...

% 第三段：深度学习编译器支持张量程序的自动生成，且伴有调优器来最大化性能...然而其普遍针对静态形状...动态形状很常见...以Qwen模型为例...GEMM算子形状有几百个（具体数字）...针对每一种形状单独调优不现实...

% 第四段：动态形状编译器（给出学术界和工业界的典型例子）被提出...但其普遍只针对单独算子来调优...未考虑整网的算子融合机会...例如CI-MI之类的...只能取得次优性能...

% 第五段：近年来...基于tile的编译器逐渐流行如Triton, TileLang, cuTile等...这类的编译器在性能和灵活性上取得了权衡..方便用户编程...我们发现其在融合算子方面有很大潜力...调优后甚至取得比厂商算子库更高的性能...但同样其调优开销过于大...不适合于动态形状...

% 第六段：从以上的分析可知，优化端到端动态形状推理存在如下挑战：1）融合方案难以确定...2）缺乏高质量编译模板...3）调优成本太大...我们提出xxx...其首先...然后...最后...

% 第七段：据我们所知...这是第一份工作xxx....本文做出如下贡献....



\begin{itemize}

\item % 我们全面分析了动态形状下端到端全局推理的性能瓶颈并给出潜在的优化机会...

\item % 在图层，我们提出xx来确认融合方案并映射到编译模板...

\item % 在算子层，我们提出xx来针对指定形状快速确定参数配置...

\item % 我们实现了xx框架...其...实验结果表明....

\end{itemize}